The \emph{Kiwitsche 2-Punkt-Verfahren} (KZV) is a practical, four-phase
procedure that uses an existing Bang-Bang (two-point) controller to identify an
FOPDT plant model and synthesise an IMC-based PID controller --- all without
requiring a separate open-loop experiment.
The method runs the Bang-Bang controller until a stationary limit cycle is
established, extracts the oscillation period and amplitude, and uses
transient-slope measurements to estimate the process gain~$K$, time
constant~$T$, and dead time~$L$.
The identified model is then fed into the Rivera--Morari IMC-PID tuning rules
to obtain the gains~$(\Kp, \Ti, \Td)$.
Simulations on a benchmark FOPDT plant show that the resulting controller
achieves \SI{5.5}{\percent} overshoot and an IAE within \SI{2}{\percent} of the
ideal IMC result, versus \SI{57.9}{\percent} overshoot and an eightfold larger
IAE for classical Ziegler--Nichols tuning.
